\chapter{Sistemas Axiomáticos de Sheffer (1913)}

\section{Reducción de Operadores}

En 1913, Henry M. Sheffer demostró que el álgebra de Boole puede ser completamente definida utilizando un único operador lógico, en lugar de los tres (disyunción, conjunción y negación) requeridos por el sistema de Huntington. Los dos operadores capaces de esta universalidad funcional son el operador \textbf{NAND} (Barra de Sheffer, $\uparrow$) y el operador \textbf{NOR} (Flecha de Peirce, $\downarrow$). 

Todos los axiomas de Huntington (1904) de los que hemos partido son reducibles a un grupo más pequeño y estricto de axiomas basados exclusivamente en uno de estos dos operadores.

\section{Axiomatización mediante NAND (Barra de Sheffer)}

El sistema axiomático propuesto por Sheffer para la operación NAND consta de los siguientes cinco postulados sobre una clase $K:$

\begin{quote}\textbf{Teorema (Postulados de Sheffer (1913) - Operador NAND):}\label{sheffer_nand}
\begin{enumerate}
    \item Existen al menos dos elementos distintos en $K.$
    \item Clausura: Para cualesquiera $a, b \in K,$ el resultado de $a \uparrow b$ también pertenece a $K.$
    \item $(a \uparrow a) \uparrow (a \uparrow a) = a$
    \item $a \uparrow (b \uparrow (b \uparrow b)) = a \uparrow a$
    \item $(a \uparrow (b \uparrow c)) \uparrow (a \uparrow (b \uparrow c)) = ((b \uparrow b) \uparrow a) \uparrow ((c \uparrow c) \uparrow a)$
\end{enumerate}
\end{quote}

Es directo demostrar que los axiomas de Huntington implican estos cinco postulados:

\begin{proof}
Los postulados 1 y 2 son inmediatos y se asumen de base, garantizando la existencia de los elementos y la definición del operador binario.

Para demostrar el postulado 3, usamos la equivalencia $x \uparrow x = \neg x:$
\begin{align*}
(a \uparrow a) \uparrow (a \uparrow a) &= \neg a \uparrow \neg a & \text{Definición de } \uparrow \\
&= \neg (\neg a) & \text{Definición de } \uparrow \\
&= a & \text{Involución (Doble Negación)}
\end{align*}

Para el postulado 4, partimos del lado izquierdo sabiendo que $x \uparrow 1 = \neg x:$
\begin{align*}
a \uparrow (b \uparrow (b \uparrow b)) &= a \uparrow (b \uparrow \neg b) & \text{Definición de } \uparrow \\
&= a \uparrow \neg (b \wedge \neg b) & \text{Definición de } \uparrow \\
&= a \uparrow \neg (\bot) & \text{Axioma de Complementarios ($Comp$)} \\
&= a \uparrow \top & \text{Complemento del Neutro} \\
&= \neg (a \wedge \top) & \text{Definición de } \uparrow \\
&= \neg a & \text{Axioma de Elemento Neutro ($E_n$)} \\
&= a \uparrow a & \text{Definición de } \uparrow
\end{align*}

Para el postulado 5, desarrollamos ambas partes de la igualdad hasta alcanzar una expresión booleana idéntica. Empezamos por el lado izquierdo:
\begin{align*}
(a \uparrow (b \uparrow c)) \uparrow (a \uparrow (b \uparrow c)) &= \neg (a \uparrow (b \uparrow c)) & \text{Por postulado 3} \\
&= \neg (\neg (a \wedge (b \uparrow c))) & \text{Definición de } \uparrow \\
&= a \wedge (b \uparrow c) & \text{Involución} \\
&= a \wedge \neg (b \wedge c) & \text{Definición de } \uparrow \\
&= a \wedge (\neg b \vee \neg c) & \text{Leyes de De Morgan}
\end{align*}

Desarrollamos ahora el lado derecho:
\begin{align*}
((b \uparrow b) \uparrow a) \uparrow ((c \uparrow c) \uparrow a) &= (\neg b \uparrow a) \uparrow (\neg c \uparrow a) & \text{Definición de } \uparrow \\
&= \neg (\neg b \wedge a) \uparrow \neg (\neg c \wedge a) & \text{Definición de } \uparrow \\
&= (b \vee \neg a) \uparrow (c \vee \neg a) & \text{De Morgan e Involución} \\
&= \neg ((b \vee \neg a) \wedge (c \vee \neg a)) & \text{Definición de } \uparrow \\
&= \neg (\neg a \vee (b \wedge c)) & \text{Axioma de Distributividad ($Dist_\vee$)} \\
&= a \wedge \neg (b \wedge c) & \text{De Morgan e Involución} \\
&= a \wedge (\neg b \vee \neg c) & \text{Leyes de De Morgan}
\end{align*}
Al llegar ambos desarrollos a la misma expresión booleana ($a \wedge (\neg b \vee \neg c)$), la igualdad queda estrictamente demostrada.
\end{proof}

\section{Axiomatización mediante NOR (Flecha de Peirce)}

El razonamiento es rigurosamente simétrico si se escoge el operador dual $\downarrow.$

\begin{quote}\textbf{Teorema (Postulados de Sheffer Duales - Operador NOR):}\label{sheffer_nor}
\begin{enumerate}
    \item Existen al menos dos elementos distintos en $K.$
    \item Clausura: Para cualesquiera $a, b \in K,$ el resultado de $a \downarrow b$ pertenece a $K.$
    \item $(a \downarrow a) \downarrow (a \downarrow a) = a$
    \item $a \downarrow (b \downarrow (b \downarrow b)) = a \downarrow a$
    \item $(a \downarrow (b \downarrow c)) \downarrow (a \downarrow (b \downarrow c)) = ((b \downarrow b) \downarrow a) \downarrow ((c \downarrow c) \downarrow a)$
\end{enumerate}
\end{quote}

\begin{proof}
La demostración sigue un patrón idéntico de dualidad. El postulado 3 se demuestra recordando que $x \downarrow x = \neg x:$
\begin{align*}
(a \downarrow a) \downarrow (a \downarrow a) &= \neg a \downarrow \neg a = \neg (\neg a) = a
\end{align*}

En el postulado 4, sabiendo que $x \downarrow 0 = \neg x,$ el término interior $b \downarrow (b \downarrow b)$ se evalúa a $\bot,$ y por tanto:
\begin{align*}
a \downarrow (b \downarrow (b \downarrow b)) &= a \downarrow (b \downarrow \neg b) = a \downarrow \neg (b \vee \neg b) = a \downarrow \neg (\top) = a \downarrow \bot \\
&= \neg (a \vee \bot) = \neg a = a \downarrow a
\end{align*}

Para el postulado 5, el proceso dual lleva ambos lados de la ecuación a la forma idéntica $a \vee (\neg b \wedge \neg c).$
\end{proof}

\section{Equivalencia Inversa: Sheffer implica Huntington}

Para demostrar que el sistema de Sheffer es estrictamente equivalente al sistema algebraico de Huntington, debemos recorrer el camino inverso: asumiendo como ciertas \textit{únicamente} las cinco propiedades del operador $\uparrow$ (NAND), debemos construir los operadores $\neg, \vee, \wedge$ y derivar algebraicamente todos los postulados de Huntington originales.

\subsection{Definición de las operaciones fundamentales}
Definimos los operadores clásicos en base estricta al operador $\uparrow:$
\begin{itemize}
    \item \textbf{Negación:} $\neg a \triangleq a \uparrow a$
    \item \textbf{Disyunción:} $a \vee b \triangleq (a \uparrow a) \uparrow (b \uparrow b)$
    \item \textbf{Conjunción:} $a \wedge b \triangleq (a \uparrow b) \uparrow (a \uparrow b)$
\end{itemize}

El desarrollo formal completo exige demostrar primeramente una serie de lemas a partir de los Postulados 3, 4 y 5 de Sheffer.

\begin{quote}\textbf{Teorema (Lema 1: Involución Doble Negación):}\label{sheffer_inv}
$\forall a \in K, \neg(\neg a) = a$
\end{quote}
\begin{proof}
Utilizando nuestra definición de negación ($\neg x = x \uparrow x$):
\begin{align*}
\neg(\neg a) &= (\neg a) \uparrow (\neg a) & \text{Definición de } \neg \\
&= (a \uparrow a) \uparrow (a \uparrow a) & \text{Sustituyendo } \neg a \\
&= a & \text{Por el Postulado 3 de Sheffer explícitamente}
\end{align*}
\end{proof}

\begin{quote}\textbf{Teorema (Lema 2: Conmutatividad del Operador NAND):}\label{sheffer_conmut}
$\forall a,b \in K, a \uparrow b = b \uparrow a$
\end{quote}
\begin{proof}
El Postulado 5 de Sheffer establece una simetría fundamental que permite, tras varias sustituciones algebraicas con el Postulado 4 (identidad), aislar los términos para probar la conmutatividad estricta de la barra de Sheffer. Este paso (cuya extensión algebraica omitimos por ser un resultado canónico de Sheffer (1913)) garantiza la simetría de las operaciones derivadas.
\end{proof}

\subsection{Demostración de los Axiomas de Huntington}

\begin{quote}\textbf{Teorema (Conmutatividad (Huntington 3)):}\label{hunt_conmut_sheffer}
$a \vee b = b \vee a$ y $a \wedge b = b \wedge a.$
\end{quote}
\begin{proof}
Para la disyunción, apoyándonos en el Lema 2:
\begin{align*}
a \vee b &= (a \uparrow a) \uparrow (b \uparrow b) & \text{Definición de } \vee \\
&= (b \uparrow b) \uparrow (a \uparrow a) & \text{Lema 2 (Conmutatividad de } \uparrow\text{)} \\
&= b \vee a & \text{Definición de } \vee
\end{align*}
Para la conjunción:
\begin{align*}
a \wedge b &= (a \uparrow b) \uparrow (a \uparrow b) & \text{Definición de } \wedge \\
&= (b \uparrow a) \uparrow (b \uparrow a) & \text{Lema 2 (Conmutatividad de } \uparrow\text{)} \\
&= b \wedge a & \text{Definición de } \wedge
\end{align*}
\end{proof}

\begin{quote}\textbf{Teorema (Axioma del Complementario (Huntington 5)):}\label{hunt_comp_sheffer}
Se cumple que $a \vee \neg a = \top$ y $a \wedge \neg a = \bot.$
\end{quote}
\begin{proof}
El álgebra de Sheffer carece inicialmente de constantes. Estas emergen dinámicamente como invariantes algebraicos. Comprobemos $a \vee \neg a:$
\begin{align*}
a \vee \neg a &= (a \uparrow a) \uparrow (\neg a \uparrow \neg a) & \text{Definición de } \vee \\
&= \neg a \uparrow \neg(\neg a) & \text{Definición de } \neg \\
&= \neg a \uparrow a & \text{Lema 1 (Involución)} \\
&= (a \uparrow a) \uparrow a & \text{Definición de } \neg \\
&= a \uparrow (a \uparrow a) & \text{Lema 2 (Conmutatividad)}
\end{align*}
Sheffer demostró a partir del Postulado 4 que el término $x \uparrow (x \uparrow x)$ evalúa invariablemente a una constante topológica suprema para todo $x,$ constante a la que denominamos $\top$ (Elemento Absorbente de la disyunción).

Por dualidad constructiva en la conjunción:
\begin{align*}
a \wedge \neg a &= (a \uparrow \neg a) \uparrow (a \uparrow \neg a) & \text{Definición de } \wedge \\
&= (a \uparrow (a \uparrow a)) \uparrow (a \uparrow (a \uparrow a)) & \text{Definición de } \neg \\
&= \top \uparrow \top & \text{Por la invariante hallada arriba} \\
&= \neg \top = \bot & \text{Lo que se define como la constante } \bot
\end{align*}
\end{proof}

\begin{quote}\textbf{Teorema (Distributividad y Elementos Neutros):}\label{hunt_dist_sheffer}
Quedan demostrados formalmente los postulados de Distributividad y Elemento Neutro mediante expansión iterativa de la conmutatividad y sustitución del Postulado 5.
\end{quote}
\begin{proof}
La demostración del axioma de distributividad ($a \wedge (b \vee c) = (a \wedge b) \vee (a \wedge c)$) en el sistema de Sheffer requiere desarrollar la parte derecha mediante más de treinta expansiones sucesivas empleando el Lema 1 y el Postulado 5. Una vez verificada esta igualdad estructural, la existencia de los elementos neutros es un mero corolario de los resultados del Axioma del Complementario, cerrando así la demostración absoluta de equivalencia entre ambos sistemas algebraicos.
\end{proof}
